\begin{frame}{Antecedentes}
\begin{itemize}
    %    \item La pandemia provocó una caída drástica en el PIB en 2020, pero para 2021 la economía ya se había recuperado.
    \item  De 2019 a 2026: 
    \begin{itemize}
        \item Los ingresos tributarios habrán aumentado menos de 1 punto del PIB...
        \item ... pero el gasto primario habrá aumentado más de 4 puntos del PIB.
        \item Las transferencias habrán aumentado casi 4 puntos del PIB, lideradas por salud y pensiones.
        \item Las últimas dos reformas tributarias no han logrado aumentar el recaudo en más de 1 punto del PIB.
    \end{itemize}
    %\item El subsidio a los combustibles continúa afectando el desempeño de las finanzas públicas.
    %\item  El próximo Gobierno se encontrará con:
    %\begin{itemize}
        %\item Altos niveles de déficit y deuda.
        %\item Intereses aún más altos por las operaciones recientes.
        %\item Liquidez escasa.
        %\item Altas reservas.
    %\end{itemize}
\end{itemize}
\end{frame}
